% ----------------------
\subsection{Section 15 (June 1829)}
% ----------------------
	\label{DC15}
	
	% -----
	\subsubsection{Order of this section relative to others around the same time}
	% -----
	
		Based on the evidence available, the JSP puts the sections from around June 1829 in the following order: \hyperref[DC14]{DC 14}, \hyperref[DC18]{DC 18}, \hyperref[DC15]{DC 15}, \hyperref[DC16]{DC 16}, and \hyperref[DC17]{DC 17}.
	
	% -----
	\subsubsection{Note on the JSP and BC versions}
	% -----
	
		The version featured in the JSP comes from BC, so for this section they are the same; the BC is therefore not included in the alternate versions below. This section doesn't appear at all in RB1.
	
	% -----
	\subsubsection{Historical Background}
	% -----
		\label{DC15-HB}
		
		% --
		\paragraph{JSP}
		% --
		
			``This revelation and the one that follows were directed to brothers John Whitmer and Peter Whitmer Jr., respectively, in early June 1829, soon after they were baptized. The two revelations are identical but were addressed to each man separately. When compiling Revelation Book 1, John Whitmer copied in the revelation addressed to himself before the revelation to Peter, implying they were received in that order.
			
			``Members of the Whitmer family became interested in JS's work as early as fall 1828. Three of the sons of Peter Whitmer Sr., John (twenty-six), David (twenty-four), and Peter Jr. (nineteen), all became `zealous friends and assistants in the work.' Peter Jr. and David were baptized sometime in June 1829 in Seneca Lake, New York, and John was probably one of those baptized the same month. On occasion, John served as a scribe to JS when Oliver Cowdery was not available to record JS's dictation of the Book of Mormon. John and Peter Jr. were `anxious to know their respective duties' and `desired with much earnestness that I [JS] should enquire of the Lord concerning them.'''
			
		% --
		\paragraph{HDDC}
		% -- 
		
			``Students of the D\&C have long noted the similarity between Sections 15 and 16. They are identical with the exception of the names of the recipients in verse 1 of each, and one additional word in verse 5 of Section 16.%
				\footnote{%
					% \label{}%
					The word is ``unto.''
					}
			The text analysis shows that these two sections were identical from 1833 to 1843, and that the additional word in verse 5 was a later insertion. These two are the only revelations of which there is a record that are word-for-word the same'' (243).
			
	% -----
	\subsubsection{Versions}
	% -----
		\label{DC15-V}
	
		\begin{tabular}{p{1in}p{2in}p{1in}p{2in}}
			JSP		& \hyperref[EMS-1828-1829-JS-JUN-1829-C]{June 1829}	& BC & Chapter 13, 33 \\
			DC 1835	& Section 40, 170									& TS & 3:20 (15 August 1842), 885 \\
			MS		& 3:10 (February 1843), 165							& DC 1844 & Section 40, 253 \\
		\end{tabular}
		
		\medskip
		
		See the \href{https://drive.google.com/file/d/1goAvNz8jS4R8slYfMG_db55Ty2z_vmgK/view?usp=sharing}{sections comparison} document for more information about the version history.

	% -----
	\subsubsection{Section 15:1} % *DC15-1 
	% -----
		\label{DC15-1}

			\footnote{%
				% \label{}%
				BC has as a heading, ``1 \textit{A Revelation given to John [Whitmer], in Fayette, New-York, June, 1829.}'' DC 1835 has, ``\textit{Revelation given to John Whitmer, jr. June, 1829.}'' DC 1844 has, `\textit{Revelation given to John Whitmer, June, 1829.}''
				}
		HEARKEN, my servant John, and listen to the words of Jesus Christ, your Lord and your Redeemer.

		\begin{framed}
		\scriptsize
			\begin{compactdesc}
				\item[JSP] HEARKEN my servant John, and listen to the words of Jesus Christ, your Lord and your Redeemer, 
				\item[DC 1835] 1 Hearken my servant John, and listen to the words of Jesus Christ, your Lord and your Redeemer, 
				\item[DC 1844] 1 Hearken my servant John, and listen to the words of Jesus Christ, your Lord and your Redeemer, 
			\end{compactdesc}
		\end{framed}

	% -----
	\subsubsection{Section 15:2} % *DC15-2 
	% -----
		\label{DC15-2}

		For behold, I speak unto you with sharpness%
			\footnote{%
				% \label{}%
				For ``sharpness,'' see \hyperref[2CORINTHIANS-13-10]{2 Corinthians 13:10}, which is the only occurrence of the word in the Bible. The Greek construction is \gr{ἀποτόμως χρήσωμαι}, which BDAG 108b says basically means \gr{ἀποτομίᾳ χρήσωμαι}, with the gloss ``that I may not have to deal sharply.'' MGLNT gives \gr{ἀποτόμως} as ``\textit{abruptly, curtly}, hence \textit{sharply, severely}.'' The word is related to ``sharper'' in \hyperref[HEBREWS-4-12]{Hebrews 4:12}, \gr{τομώτερος}. Speaking of the noun \gr{ἀποτομία}, MM 71b notes a papyrus use and says, ``Counsel is pleading a native statute, admittedly harsh, which he claims was enforced rigidly: the word does not suggest straining a statute, but simply exacting its provisions to the full.'' So think maybe also of the exactness of the commandments and how for God there can be no deviation from them.
				}
		and with power, for mine arm is over all the earth.

		\begin{framed}
		\scriptsize
			\begin{compactdesc}
				\item[JSP] for behold I speak unto you with sharpness and with power, for mine arm is over all the earth, 
				\item[DC 1835] for behold I speak unto you with sharpness and with power, for mine arm is over all the earth, 
				\item[DC 1844] for behold I speak unto you with sharpness and with power, for mine arm is over all the earth, 
			\end{compactdesc}
		\end{framed}

	% -----
	\subsubsection{Section 15:3} % *DC15-3 
	% -----
		\label{DC15-3}

		And I will tell you that which no man knoweth save me and thee alone---

		\begin{framed}
		\scriptsize
			\begin{compactdesc}
				\item[JSP] and I will tell you that which no man knoweth save me and thee alone: 
				\item[DC 1835] and I will tell you that which no man knoweth save me and thee alone: 
				\item[DC 1844] and I will tell you that which no man knoweth save me and thee alone, 
			\end{compactdesc}
		\end{framed}

	% -----
	\subsubsection{Section 15:4} % *DC15-4 
	% -----
		\label{DC15-4}

		For many times you have desired of me to know that which would be of the most worth unto you.

		\begin{framed}
		\scriptsize
			\begin{compactdesc}
				\item[JSP] for many times you have desired of me to know that which would be of the most worth unto you.
				\item[DC 1835] for many times you have desired of me to know that which would be of the most worth unto you.
				\item[DC 1844] for many times have you desired of me to know that which would be of the most worth unto you.
			\end{compactdesc}
		\end{framed}

	% -----
	\subsubsection{Section 15:5} % *DC15-5 
	% -----
		\label{DC15-5}

		Behold, blessed are you for this thing, and for speaking my words which I have given you according to my commandments.

		\begin{framed}
		\scriptsize
			\begin{compactdesc}
				\item[JSP] 2 Behold, blessed are you for this thing, and for speaking my words which I have given you, according to my commandments:
				\item[DC 1835] 2 Behold, blessed are you for this thing, and for speaking my words which I have given you, according to my commandments:
				\item[DC 1844] 2 Behold, blessed are you for this thing, and for speaking my words which I have given you, according to my commandments.
			\end{compactdesc}
		\end{framed}

	% -----
	\subsubsection{Section 15:6} % *DC15-6 
	% -----
		\label{DC15-6}

		And now, behold, I say unto you, that the thing which will be of the most worth unto you will be to declare repentance unto this people, that you may bring souls unto me, that you may rest with them in the kingdom of my Father. Amen.

		\begin{framed}
		\scriptsize
			\begin{compactdesc}
				\item[JSP] 3 And now behold I say unto you, that the thing which will be of the most worth unto you, will be to declare repentance unto this people, that you may bring souls unto me, that you may rest with them in the kingdom of my Father. Amen.
				\item[DC 1835] 3 And now behold I say unto you, that the thing which will be of the most worth unto you, will be to declare repentance unto this people, that you may bring souls unto me, that you may rest with them in the kingdom of my Father. Amen.
				\item[DC 1844] 3 And now behold I say unto you, that the thing which will be of the most worth unto you, will be to declare repentance unto this people, that you may bring souls unto me, that you may rest with them in the kingdom of my Father: Amen.
			\end{compactdesc}
		\end{framed}